\prefacesection{Preface}

The work presented in this dissertation grew out of a deceptively simple question: why do we remember some moments well, and others not at all --- even when the to-be-remembered content is held constant? Across my years of training in the Stanford Memory Lab, I came to appreciate that a meaningful part of the answer lies not in the content of memory itself, but in the seconds just before we attempt to access it. Whether attention is focused or drifting, whether arousal sits near optimal or wanders toward its extremes, shapes the probability that an otherwise well-encoded experience can be brought back to mind.

This dissertation pursues that intuition along four intersecting threads --- theory, measurement, conceptual framework, and causal experiment --- presented as a collection of three first-author papers and one previously unpublished framework chapter, bookended by an introductory framing chapter and a concluding synthesis.

Chapter 2 is reproduced from a review article published in \textit{Current Directions in Psychological Science} and co-authored with Haopei Yang, Alice M. Xue, and my advisor, Anthony D. Wagner. The article surveys recent advances in methods and theory that illuminate how attention shapes learning and remembering, and it lays out the conceptual scaffolding that motivates the empirical and methodological work that follows.

Chapter 3 is reproduced from a preprint co-authored with Haopei Yang, Alice M. Xue, and Mingjian He, introducing \texttt{eyeris}, an R package for transparent and reproducible pupillometry preprocessing. The package is offered both as a methodological prerequisite for the closed-loop work that follows and, to the broader research community, as a standalone contribution toward open and reproducible psychophysiological science.

Chapter 4 is reproduced from an unpublished paper drafted early in my doctoral training. The chapter lays out a conceptual framework and taxonomy for closed-loop pupillometry interventions on attention and memory. Although it never matured into a published manuscript, it is included here because it foreshadows the design choices that defined the empirical paradigm developed in the chapter that follows, and because this dissertation is the most fitting outlet in which to preserve the conceptual lineage of that work.

Chapter 5 describes PRISM, a closed-loop pupillometry paradigm that uses real-time readouts of arousal and attention to causally test the readiness-to-remember framework introduced in Chapter 2, that depends on the preprocessing infrastructure developed in Chapter 3, and that constitutes one empirical realization of the framework outlined in Chapter 4. This chapter is in preparation for submission as a co-authored manuscript.

The introductory (Chapter 1) and concluding (Chapter 6) chapters were written specifically for this dissertation and are not part of any prior publication. The introduction situates the four substantive chapters within a broader theoretical landscape, and the conclusion draws their findings together to articulate open questions and promising future directions.
